\documentclass[tikz, border=10pt]{standalone}
\usepackage{pgfplots}
\pgfplotsset{compat=1.18}
\usetikzlibrary{patterns} % Librería necesaria para el patrón de rayas

\begin{document}
	\definecolor{maincolor1}{RGB}{1,82,142}
	\definecolor{maincolor2}{RGB}{21,182,184}
	\definecolor{mygrey}{rgb}{0.41, 0.41, 0.41}
	\definecolor{discardred}{RGB}{210, 60, 60}
	\definecolor{reusecolor}{RGB}{230, 150, 30}

	\begin{tikzpicture}
		\begin{axis}[
			title={$f(x) = x^2 - 3.2x + 3$ \quad (It. 3)},
			xlabel={$x$},
			ylabel={$f(x)$},
			xmin=-0.2, xmax=3.2,
			ymin=-0.1, ymax=3.3,
			axis lines=left,
			grid style={line width=.1pt, draw=gray!20},
			major grid style={line width=.2pt, draw=gray!50},
			minor tick num=1,
			legend pos=north west,
			legend style={font=\small, fill=white, fill opacity=0.9, draw opacity=1},
			width=12cm,
			height=8cm
			]

			% Región descartada [0, 1.1459] con rayas diagonales rojas
			\addplot[
			draw=none,
			pattern=north east lines, 
			pattern color=discardred
			] coordinates {
				(0, -0.1) (0, 3.3) (1.1459, 3.3) (1.1459, -0.1)
			};

			% Región descartada [1.8541, 3] con rayas diagonales rojas
			\addplot[
			draw=none,
			pattern=north east lines, 
			pattern color=discardred
			] coordinates {
				(1.8541, -0.1) (1.8541, 3.3) (3, 3.3) (3, -0.1)
			};

			% Nuevo intervalo de incertidumbre aceptado [1.1459, 1.8541] (Fondo gris)
			\addplot[
			draw=none,
			fill=mygrey, 
			opacity=0.1
			] coordinates {
				(1.1459, -0.1) (1.1459, 3.3) (1.8541, 3.3) (1.8541, -0.1)
			};

			% Curva continua f(x)
			\addplot [
			domain=0:3,
			samples=100,
			color=maincolor1,
			thick
			] {x^2 - 3.2*x + 3};

			% Líneas proyectadas para los puntos de prueba x1 = 1.5836 y x2 = 1.8541
			\draw[gray!50, dashed] (axis cs:1.5836, -0.1) -- (axis cs:1.5836, 0.4403);
			\draw[gray!50, dashed] (axis cs:1.8541, -0.1) -- (axis cs:1.8541, 0.5046);

			% Puntos de prueba evaluados en esta iteración
			\addplot [
			only marks, 
			mark=*, 
			mark size=3pt, 
			color=maincolor2
			] coordinates {
				(1.5836, 0.4403)
			};

			% Punto reutilizado de la iteración anterior (no se vuelve a evaluar)
			\addplot [
			only marks,
			mark=*,
			mark size=3pt,
			color=reusecolor
			] coordinates {(1.8541, 0.5046)};

			% Mínimo analítico exacto (1.6, 0.44)
			\addplot [
			only marks,
			mark=diamond*,
			mark size=4pt,
			color=maincolor2,
			thick
			] coordinates {(1.6, 0.44)};

			% Etiquetas en el eje X para x1 y x2
			\node[below, font=\footnotesize, maincolor1] at (axis cs:1.5836, -0.1) {$x_1$};
			\node[below, font=\footnotesize, maincolor1] at (axis cs:1.8541, -0.1) {$x_2$};

		\end{axis}
	\end{tikzpicture}
\end{document}
